\chapter{Results}
We begin by discussing the constraints for UV fields that would produce the operator $\mathcal{O}_{le}$. From the LEFT bounds in Tab.~\ref{tab:leftbounds} and from the matching condition in Eq.~\ref{eq:specMatchVLR2112} we know that the real and imaginary part of its coefficient lie in this $95\%$-confidence region:
\begin{equation}
    \text{Re}[C_{le}^{2112}]\text{ in }[-0.021,+0.018]\text{ and Im}[C_{le}^{2112}]\text{ in }[-0.017,+0.023]
\end{equation}
The UV matching of this operator was cited in Eq.~\ref{eq:UVMatchCle} and for $C_{le}^{2112}$ amounts to:
\begin{equation}
     C_{le}^{2112} =
     \frac{\mathcal{D}_{{e^\dagger eV_{1}}}^{{12p_1}} \mathcal{D}_{{ll^\dagger V_{1}}}^{{12p_1}}}{{M_{V_{1}}^2}}
    -\frac{\mathcal{D}_{{e^\dagger lS_{4}^\dagger }}^{{11p_1}} \mathcal{D}_{{e^\dagger lS_{4}^\dagger }}^{{22p_1*}}}{2 {M_{S_{4}}^2}}
    +\frac{\mathcal{D}_{{e^\dagger l^\dagger V_{3}^\dagger }}^{{12p_1}} \mathcal{D}_{{e^\dagger l^\dagger V_{3}^\dagger }}^{{21p_1*}}}{{M_{V_{3}}^2}}
\end{equation}
If we assume single-mediator dominance, a single heavy copy accordingly, and a unit real coupling product, the appropriate limits for the masses are:
\begin{align}
    M_{V_1},M_{V_3}&>\frac{v}{\sqrt{0.018}}\simeq1.84\text{ TeV}\\
    M_{S_4}&>\frac{v}{\sqrt{2(0.021)}}\simeq1.2\text{ TeV}
\end{align}
For $S_4$ we used the lower end of the $95\%$ CL region and for $V_1,V_3$ the upper end due to the sign in their matching condition. The SMEFT RGEs only change the Wilson coefficient by about $1\%$ as we see from Tab.~\ref{tab:smeft-rge-primary}.
\section{Heavy Scalar \texorpdfstring{$\mathbf{S_4}$}{S4}}
It's time to \textit{change gears} and reverse the direction: We take a look at which SMEFT operators $S_4$ produces and if they would show up in muon decay. From the dictionary we can look up the top-down translation of $S_4$:
\begin{equation}
    S_4\to \op_H,\op_{eH},\op_{uH},\op_{dH},\op_{le},\op^{(1)}_{qu},\op^{(8)}_{qu},\op^{(1)}_{qd},\op^{(8)}_{qd},
    \op_{leqd},\op_{quqd}^{(1)},\op^{(1)}_{lequ}
\end{equation}
But we do can not really make any educated assumptions on the hypothetical coupling of $S_4$ to the quark sector or the Higgs. Let us assume those other couplings are zero -- can we infer something nevertheless? \par
The Lagrangian term of $S_4$ relevant for its matching conditions is this:
\begin{align}
\begin{autobreak}
{\Delta\mathcal{L}_{S_4}} =
- 2 \mathcal{D}_{{d^\dagger QS_{4}^\dagger }}^{{rsp}} [(\bar{d}_{{r}})^{{a}}({q}_{{s}})_{{ai}}]({S}_{{4p}}^{\dagger })^{{i}} 
- \mathcal{D}_{{e^\dagger lS_{4}^\dagger }}^{{rsp}} [\bar{e}_{{r}}({l}_{{s}})_{{i}}]({S}_{{4p}}^{\dagger })^{{i}} 
- 2 \mathcal{D}_{{QS_{4}u^\dagger }}^{{rps}}\epsilon ^{{ji}} [(\bar{u}_{{s}})^{{a}}({q}_{{r}})_{{aj}}]({S}_{{4p}})_{{i}}  
- \mathcal{D}_{{HH^\dagger H^\dagger S_{4}(1)}}^{{p}}\delta _{{k}}^{{i}} \delta _{{l}}^{{j}}{H}_{{j}}{H}^{{\dagger k}}{H}^{{\dagger l}}({S}_{{4p}})_{{i}} 
- \mathcal{D}_{{HH^\dagger H^\dagger S_{4}(2)}}^{{p}}\delta _{{l}}^{{i}} \delta _{{k}}^{{j}}{H}_{{j}}{H}^{{\dagger k}}{H}^{{\dagger l}}({S}_{{4p}})_{{i}}
\end{autobreak}\label{eq:UVLagS4}
\end{align}
So the matching condition of $C_{le}$ to $S_4$ in $prst$ convention is the following:
\begin{equation}
 {C_{le,prst}} = -\frac{\mathcal{D}_{{e^\dagger lS_{4}^\dagger }}^{{sr}} \mathcal{D}_{{e^\dagger lS_{4}^\dagger }}^{{tp*}}}{2 {M_{S_{4}}^2}} \label{eq:UVMatchCle}
\end{equation}
We get for $C_{le,2112}$ the following matching:
\begin{equation}
{C_{le,2112}} = -\frac{\mathcal{D}_{{e^\dagger lS_{4}^\dagger }}^{{11a}} \mathcal{D}_{{e^\dagger lS_{4}^\dagger }}^{{22a*}}}{2 {M_{S_{4}}^2}}=-\frac{\lambda_e\lambda_\mu^*}{2M_{S_4}^2}
\end{equation}
We thus can infer:
\begin{equation}
{|C_{le,1111}}| = -\frac{\lambda_e^2}{2M_{S_4}^2}\qquad{|C_{le,2222}|} = -\frac{\lambda_\mu^2}{2M_{S_4}^2}\qquad{|C_{le,1221}|} = -\frac{\lambda_e^*\lambda_\mu}{2M_{S_4}^2}
\end{equation}
Together with our assumptions $\lambda_e\simeq1$ and $\lambda_\mu\simeq1$ this translates to:
\begin{equation}
    |C_{le}^{1111}| > 1.2\text{ TeV}\qquad{|C_{le}^{2222}|} > 1.2\text{ TeV} \qquad{|C_{le}^{2112}|} > 1.2\text{ TeV}
\end{equation}
For $\op_{le,1111}=(\bar l_e\gamma_\mu l_e)(\bar e_R\gamma^\mu e_R)$ probably Bhabha scattering $e^+e^-\to e^+e^-$ provides the most direct probe. Here, $C_{le,1111}$ interferes with the SM photon and Z-exchange amplitudes. A recent scalar-doublet analysis using the LEP differential distributions quotes \cite{erdelyi2025probing}:
\begin{equation}
    \frac{M_{S_4}}{|\lambda_e|}>2.7\text{ TeV}\qquad (95\%\text{ CL})
\end{equation}
At $M_{S_4}=2.7\text{ TeV}$, the unit-coupling prediction is already:
\begin{equation}
    |v^2C_{le}^{2112}|=\frac{v^2}{2M_{S_4}^2}\simeq 0.00416
\end{equation}
This is well inside the Michel parameter interval. However, muon decay still gives complimentary sensitivity to the complex phase. For the coefficient $C_{le,2222}$ a future high-energy muon collider could probe the muonic analogue of Bhabha scattering, $\mu^+\mu\to\mu^+\mu^-$ and significantly extend current sensitivity to muon contact interactions \cite{glioti2025flavor}.\par 
As $S_4$ is an electroweak scalar doublet, it contains charged and neutral components. The gauge interactions are fixed by its representation, independently of $\lambda_e$ or $\lambda_\mu$, and the components could thus be produced through electroweak processes involving $\gamma,Z,W^\pm$. The expected decay would be then:
\begin{align}
    S^0_4&\to e^+e^-\qquad S^0_4\to \mu^+\mu^-\\
    S^\pm_4&\to e^\pm\nu_e\qquad S^\pm_4\to\mu^\pm\nu_\mu
\end{align}
It has been shown that a charged scalar decaying into electrons or muons would provide an especially clean collider signature \cite{alcaide2020lhc}.\par
Falkowski \& Mimouni discuss in \cite{falkowski2016model} the constraints on SMEFT couplings from Michel parameter data as well and remark, like us, that at linear order the best bound comes from $\eta$ and $\beta'/A$ on the real and imaginary part of $C_{le}^{2112}$ ($[c_{le}]_{1221}$ in their notation) respectively. They refrain from giving an explicit bound on $S_4$ (\enquote{inert Higgs}) exactly because LEP data provides the stronger bound. 

\section{Correlated mediator \texorpdfstring{$\mathbf{V_1}$}{V1}}
The heavy mediator $V_1$ produces multiple LEFT coefficients at the scale of muon decay and we use it to obtain a mass bound from a correlated fit. Let us first take a look at its Lagrangian:
\begin{align}
\begin{autobreak}
    {\Delta\mathcal{L}_{V_4}=}

- 2 \mathcal{D}_{{d^\dagger dV_{1}}}^{{rsp}} [(\bar{d}_{{r}})^{{a}}\gamma _{\mu }({d}_{{s}})_{{a}}]{V}_{{1p}}^{\mu } 
 - \mathcal{D}_{{e^\dagger eV_{1}}}^{{rsp}}1[\bar{e}_{{r}}\gamma _{\mu }{e}_{{s}}]{V}_{{1p}}^{\mu } 
 - 2 \mathcal{D}_{{HH^\dagger V_{1}D}}^{{p}} [D_{\mu }{H}_{{i}}]{H}^{{\dagger i}}{V}_{{1p}}^{\mu } 
 + \mathcal{D}_{{ll^\dagger V_{1}}}^{{rsp}} [(\bar{l}_{{s}})^{{i}}\gamma _{\mu }({l}_{{r}})_{{i}}]{V}_{{1p}}^{\mu } 
 + 2 \mathcal{D}_{{QQ^\dagger V_{1}}}^{{rsp}} [(\bar{q}_{{s}})^{{ai}}\gamma _{\mu }({q}_{{r}})_{{ai}}]{V}_{{1p}}^{\mu } 
 - 2 \mathcal{D}_{{u^\dagger uV_{1}}}^{{rsp}} [(\bar{u}_{{r}})^{{a}}\gamma _{\mu }({u}_{{s}})_{{a}}]{V}_{{1p}}^{\mu } 
\end{autobreak}\label{eq:V1Lag}
\end{align}
We choose real unit couplings $\lambda_l=\lambda_e=h=1$ and define:
\begin{align}
    s\equiv \frac{v}{M_V}\qquad v=0.246\text{ TeV}\\
    x_{ab}\equiv v^2 L_{\nu e,ab12}^{V,LL}\qquad y_{ab}\equiv v^2 L_{\nu e,ab12}^{V,LR}
\end{align}
The nonzero coefficients are then:
\begin{equation}
    \begin{aligned}
         (a,b) &\quad x_{ab} & y_{ab}  \\[2mm]
         (2,1) &\quad -s^2   & s^2     \\
         (1,2) &\quad -s^2   & s^2     \\
         (1,1),(2,2),(3,3) &\quad s^2   & -s^2 
    \end{aligned}
\end{equation}
If we substitute this into our full-fermion formulas from chapter 3 we receive:
\begin{align}
    R_G(s)\equiv\frac{G_F^2}{G_{F,\rm SM}}&=1+s^2+\frac{5}{2}s^4\\
    \rho(s)=\delta(s)&=\frac{3}{4}\\
    \xi(s)=\xi'(s)&=1-\frac{5}{2}s^4\\
    \xi''(s)&=1\\
    \eta(s)=-\eta''(s)&=\frac{1}{2}s^2+\frac{3}{4}^4\\
    \frac{\alpha'}{A}=\frac{\beta'}{A}&=0
\end{align}
Define a total chi-squared over our five experimental blocks (TWIST, Danneberg, $\xi$, $\xi'$ and the electroweak input scheme for $G_{F,\rm SM}$):
\begin{equation}
    \chi^2(s)=\chi^2_T(s)+\chi^2_D(s)+\left(\frac{\xi'(s)-1.00}{0.04}\right)^2+\left(\frac{\xi''(s)-0.98}{0.04}\right)^2+\chi_{G_F}^2(s)
\end{equation}
And we obtain:
\begin{equation}
    M_{V_1}>\frac{v}{s_{95}}=\frac{0.246\text{ TeV}}{0.0595}=4.13\text{ TeV}
\end{equation}
This is consistent with our first naïve result in Eq.~\ref{eq:naiveresult} and moreover with the findings of Falkowski \& Mimouni in Sec.~5.2 of \cite{falkowski2016model} where they report $M_V/\lambda_{l,e}>3.3\text{ to }10\text{ TeV}$ depending on the ratio of left- and right-handed couplings.

\subsubsection{Input-parameter effects of $\mathbf{C_{HD}}$}
\label{sec:v1-chd-input}

The neutral vector singlet $V_1\sim(\mathbf 1,\mathbf 1)_0$ may couple to
the Hermitian Higgs current according to:
\begin{equation}
    \Delta\mathcal L_{V_1}
    \supset
    -\lambda_h V_{1,\mu}J_H^\mu
    \qquad
    J_H^\mu=H^\dagger i\overleftrightarrow D^\mu H
\end{equation}
Integrating out $V_1$ generates, among other operators:
\begin{equation}
    C_{HD}=-\frac{2\lambda_h^2}{M_{V_1}^2}
    \label{eq:v1-chd-matching-short}
\end{equation}
This coefficient does not induce a four-lepton contact interaction or
modify the charged-current $Wl\nu$ vertex directly.  Its effect on
muon decay instead arises through the relations between the Lagrangian
parameters and the measured electroweak input observables.  Writing the
Higgs doublet after electroweak symmetry breaking as
\begin{equation}
    H=\frac{1}{\sqrt2}
      \begin{pmatrix}0\\v+h\end{pmatrix},
\end{equation}
one finds, to linear order in $C_{HD}$:
\begin{equation}
    m_W^2=\frac{g_2^2v^2}{4}
    \qquad
    m_Z^2=\frac{(g_1^2+g_2^2)v^2}{4}
    \left(1+\frac12v^2C_{HD}\right)
    \label{eq:chd-input-masses-short}
\end{equation}
Thus $C_{HD}$ changes the values of $g_1$, $g_2$ and $v$ inferred from a
chosen set of electroweak inputs.  In the
$\{\widehat m_W,\widehat m_Z,\widehat\alpha\}$ scheme used here, keeping
the measured inputs fixed gives:
\begin{equation}
    \frac{\delta v}{v}
    =-\frac{c_W^2}{4s_W^2}v^2C_{HD}
    \qquad
    \frac{\delta G_F^{\mathrm{pred}}}{G_F^{\mathrm{pred}}}
    =\frac{c_W^2}{2s_W^2}v^2C_{HD}
    =-\frac{c_W^2}{s_W^2}
      \frac{\lambda_h^2v^2}{M_{V_1}^2}
    \label{eq:chd-input-gf-short}
\end{equation}
Since the decay rate scales as $\Gamma_\mu\propto G_F^2$, this produces
an overall normalization shift:
\begin{equation}
    \left.\frac{\delta\Gamma_\mu}{\Gamma_\mu}\right|_{C_{HD}}
    =-2\frac{c_W^2}{s_W^2}
      \frac{\lambda_h^2v^2}{M_{V_1}^2}.
\end{equation}
It does not, at this order, change the normalized Michel-spectrum shape
parameters.  If $G_F$ is itself chosen as an electroweak input, the same
effect is absorbed into the inferred value of $v$ and instead appears in
the predictions for other electroweak observables.

A consistent treatment would therefore require the simultaneous
redefinition of the electroweak input parameters and the inclusion of
electroweak precision observables together with the remaining operators
generated by $V_1$.  Such a global analysis lies beyond the scope of this
work and is not pursued further.  

\section{Dimension seven results}\label{sec:results-d7}
The constraints from the dimension-seven operators are considerably lower than the ones from the dimension-six SMEFT -- this is because their leading contribution to the Michel observables scale as.
\begin{equation}
\Delta\mathcal{O}\sim |z_7|^2
\sim\left(\frac{v}{\Lambda}\right)^6
\end{equation}
In \cite{fridell2024probing} the authors constrain the operator $\op_{\bar elllH}$ via non-standard muon decay: $L^{S,LL}_{\nu e,2112}$ leads to $\Lambda_\text{NP}\approx 0.4\text{ TeV}$
and $L^{S,LL}_{\nu e,2112}$ to $\Lambda_\text{NP}\approx 0.2\text{ TeV}$ under the assumption of single LEFT operator dominance to $\mu^+\to e^+ +\bar\nu_e+\bar\nu_\mu$. The results agree at the expected order of magnitude -- our somewhat higher coefficient-scale sensitivity follows from a tighter Michel-parameter constraint and from using projected symmetric or mixed flavor directions rather than the single raw coefficient employed in the earlier estimate.
